\documentclass[11pt,a4paper]{article}
\usepackage[italian]{babel}
\usepackage[utf8]{inputenc}
\usepackage{amsmath}
\usepackage{amssymb}
\usepackage{amsthm}  % Per l'ambiente example
\usepackage{graphicx}
\usepackage{booktabs}
\usepackage{array}
\usepackage{multirow}
\usepackage{listings}
\usepackage{xcolor}
\usepackage{geometry}
\usepackage{enumitem}
\usepackage{float}
\usepackage{caption}
\usepackage{subcaption}
\usepackage{textcomp}  % Per caratteri speciali

% Definisci l'ambiente example (opzionale)
\newtheorem{example}{Esempio}

% Imposta margini più larghi per evitare Overfull \hbox
\geometry{a4paper, left=2.5cm, right=2.5cm, top=2cm, bottom=2cm}

% Stile per listings (riduci la dimensione del font)
\definecolor{codegreen}{rgb}{0,0.6,0}
\definecolor{codegray}{rgb}{0.5,0.5,0.5}
\definecolor{codepurple}{rgb}{0.58,0,0.82}
\definecolor{backcolour}{rgb}{0.95,0.95,0.92}

\lstdefinestyle{mystyle}{
    backgroundcolor=\color{backcolour},
    basicstyle=\ttfamily\tiny,  % Dimensione ridotta
    breaklines=true,
    frame=single,
    rulecolor=\color{gray!30}
}
\lstset{style=mystyle}

% Permetti spazi flessibili
\sloppy

\title{Spiegazione LiveDataset}
\author{}
\date{}

\begin{document}

\maketitle

\section{Obiettivo del progetto}
Il progetto contenuto nel notebook \texttt{liverDataSet(1).ipynb} affronta un problema di \textit{machine learning supervisionato} di tipo \textit{classificazione binaria}.

L'obiettivo è costruire un modello capace di prevedere se un paziente appartiene a una delle due classi:
\begin{itemize}
    \item \texttt{Liver Disease}
    \item \texttt{No Liver Disease}
\end{itemize}

cioè se il paziente presenta oppure no una malattia epatica, sulla base di alcune variabili cliniche come età, sesso, valori di bilirubina, enzimi epatici e proteine nel sangue.

In termini di \textit{machine learning}, abbiamo:
\begin{align*}
    \text{Input}  &= \text{dati clinici del paziente} \\
    \text{Output} &= \text{classe del paziente}
\end{align*}

La domanda principale del progetto è:
\begin{quote}
    dati certi valori medici, è possibile addestrare un modello che distingua in modo abbastanza affidabile i pazienti con malattia epatica da quelli senza malattia epatica?
\end{quote}

\section{Dataset utilizzato}
Il dataset viene scaricato da Kaggle tramite \texttt{kagglehub}:

\begin{lstlisting}[language=Python]
import kagglehub

path = kagglehub.dataset_download(
    "shauryasrivastava01/liver-patient-dataset"
)
\end{lstlisting}

Successivamente viene letto il file CSV:

\begin{lstlisting}[language=Python]
path = "Datasets/liver_patient_dataset.csv"
df = pd.read_csv(path)
\end{lstlisting}

Il dataset contiene \textbf{583 righe} e \textbf{11 colonne}. Ogni riga rappresenta un paziente.

\begin{table}[H]
\centering
\caption{Descrizione delle colonne del dataset}
\label{tab:dataset_columns}
\begin{tabular}{>{\raggedright\arraybackslash}p{3cm}>{\raggedright\arraybackslash}p{6cm}}
\toprule
\textbf{Colonna} & \textbf{Significato} \\
\midrule
\texttt{Age} & età del paziente \\
\texttt{Gender} & sesso del paziente \\
\texttt{TB} & Total Bilirubin \\
\texttt{DB} & Direct Bilirubin \\
\texttt{Alkphos} & Alkaline Phosphotase \\
\texttt{Sgpt} & enzima epatico SGPT / ALT \\
\texttt{Sgot} & enzima epatico SGOT / AST \\
\texttt{TP} & Total Proteins \\
\texttt{ALB} & Albumin \\
\texttt{A/G Ratio} & rapporto albumina/globulina \\
\texttt{Selector} & variabile target \\
\bottomrule
\end{tabular}
\end{table}

La variabile più importante per il progetto è \texttt{Selector}, perché è il \textit{target}, cioè la classe da predire.

\section{Variabile target: \texttt{Selector}}
La colonna \texttt{Selector} contiene due valori testuali:
\begin{itemize}
    \item \texttt{Liver Disease}
    \item \texttt{No Liver Disease}
\end{itemize}

Nel notebook viene visualizzata la distribuzione delle classi:

\begin{lstlisting}[language=Python]
conteggi = df["Selector"].value_counts()

plt.bar(conteggi.index, conteggi.values)
plt.xlabel("Selector")
plt.ylabel("Frequenza")
plt.title('Distribuzione di "Selector"')
plt.show()
\end{lstlisting}

Questa parte serve per capire se il dataset è \textit{bilanciato} oppure \textit{sbilanciato}.

Un dataset è \textbf{bilanciato} quando le classi hanno circa lo stesso numero di esempi. Esempio ideale:
\begin{align*}
    \text{Liver Disease}     &= 300 \\
    \text{No Liver Disease}  &= 300
\end{align*}

Un dataset è \textbf{sbilanciato} quando una classe è molto più frequente dell'altra. In questo caso il dataset è sbilanciato: i pazienti con \texttt{Liver Disease} sono molti di più dei pazienti con \texttt{No Liver Disease}.

Questo è un punto fondamentale del progetto, perché uno sbilanciamento può ingannare il modello.

\section{Perché lo sbilanciamento è un problema?}
Supponiamo di avere questo dataset:
\begin{align*}
    90 \text{ pazienti con malattia} \\
    10 \text{ pazienti senza malattia}
\end{align*}

Un modello molto stupido potrebbe predire sempre:
\begin{quote}
    \texttt{Liver Disease}
\end{quote}

In quel caso avrebbe:
\begin{align*}
    90\% \text{ di accuracy}
\end{align*}

ma sarebbe un pessimo modello, perché non riconoscerebbe mai la classe minoritaria. Quindi \textbf{l'accuracy da sola non basta}.

Nel progetto, infatti, vengono usate anche metriche come:
\begin{itemize}
    \item precision
    \item recall
    \item f1-score
    \item balanced accuracy
    \item confusion matrix
\end{itemize}

Queste metriche sono molto più informative quando il dataset è sbilanciato.

\section{Separazione tra feature e target}
Nel notebook vengono separate le variabili di input dalla variabile target:

\begin{lstlisting}[language=Python]
feature_names = set(df.columns) - set(["Selector"])
target_name = "Selector"
categorical_features = set(["Gender"])
continuos_features = feature_names - set(["Gender"])
\end{lstlisting}

Qui si dice:
\begin{quote}
    Tutte le colonne tranne \texttt{Selector} sono feature. \texttt{Selector} è il target. \texttt{Gender} è categorica. Le altre variabili sono continue/numeriche.
\end{quote}

Quindi:
\begin{align*}
    X &= \text{df senza Selector} \\
    y &= \text{df["Selector"]}
\end{align*}

Dove:
\begin{align*}
    X &= \text{dati usati per predire} \\
    y &= \text{valore da predire}
\end{align*}

\begin{table}[H]
\centering
\caption{Esempio di riga del dataset}
\label{tab:dataset_example}
\begin{tabular}{lcccccccccl}
\toprule
\textbf{Age} & \textbf{Gender} & \textbf{TB} & \textbf{DB} & \textbf{Alkphos} & \textbf{Sgpt} & \textbf{Sgot} & \textbf{TP} & \textbf{ALB} & \textbf{A/G Ratio} & \textbf{Selector} \\
\midrule
65 & Female & 0.7 & 0.1 & 187 & 16 & 18 & 6.8 & 3.3 & 0.90 & Liver Disease \\
\bottomrule
\end{tabular}
\end{table}

Il modello vede:
\begin{align*}
    \text{Age} &= 65 \\
    \text{Gender} &= \text{Female} \\
    \text{TB} &= 0.7 \\
    \text{DB} &= 0.1 \\
    &\vdots
\end{align*}

e deve predire:
\begin{align*}
    \text{Selector} = \text{Liver Disease}
\end{align*}

\section{Codifica delle variabili categoriche}
I modelli di scikit-learn non lavorano direttamente con stringhe come:
\begin{itemize}
    \item \texttt{Female}
    \item \texttt{Male}
    \item \texttt{Liver Disease}
    \item \texttt{No Liver Disease}
\end{itemize}

Per questo nel notebook si crea una copia codificata del dataset:

\begin{lstlisting}[language=Python]
df_encoded = df.copy()

df_encoded["Gender"] = df_encoded["Gender"].astype("category").cat.codes
df_encoded["Selector"] = df_encoded["Selector"].astype("category").cat.codes
\end{lstlisting}

Dopo questa trasformazione:
\begin{align*}
    \text{Female} &\rightarrow 0 \\
    \text{Male} &\rightarrow 1 \\
    \text{Liver Disease} &\rightarrow 0 \\
    \text{No Liver Disease} &\rightarrow 1
\end{align*}

Questa cosa è molto importante. Nel notebook, quando si parla di \texttt{class\_1}, si sta parlando di:
\begin{quote}
    \texttt{No Liver Disease}
\end{quote}

non di:
\begin{quote}
    \texttt{Liver Disease}
\end{quote}

Quindi se il codice ottimizza \texttt{f1\_class\_1}, sta cercando di migliorare il riconoscimento dei pazienti \textbf{senza} malattia epatica.

\section{Divisione train/test}
Il dataset viene diviso in \textit{training set} e \textit{test set}:

\begin{lstlisting}[language=Python]
from sklearn.model_selection import train_test_split

X_train, X_test, y_train, y_test = train_test_split(
    df_encoded[feature_names],
    df_encoded[target_name],
    test_size=0.3,
    random_state=42,
    stratify=y
)
\end{lstlisting}

La divisione è:
\begin{align*}
    70\% &\rightarrow \text{training set} \\
    30\% &\rightarrow \text{test set}
\end{align*}

Il \textit{training set} serve per addestrare il modello. Il \textit{test set} serve per valutare il modello su dati mai visti.

Il parametro più importante qui è:
\begin{lstlisting}[language=Python]
stratify=y
\end{lstlisting}

Questo serve a mantenere nel train e nel test una distribuzione simile delle classi. Per esempio, se nel dataset completo circa il 70\% dei pazienti è nella classe 0 e il 30\% nella classe 1, anche il training set e il test set manterranno più o meno quelle proporzioni. Questo è utile soprattutto quando il dataset è sbilanciato.

\section{Primo modello: Decision Tree senza vincoli}
Il primo modello addestrato è un \textit{albero decisionale}:

\begin{lstlisting}[language=Python]
from sklearn.tree import DecisionTreeClassifier

model = DecisionTreeClassifier(random_state=42)
model = model.fit(X_train, y_train)
\end{lstlisting}

Un albero decisionale è un modello che prende decisioni tramite una sequenza di domande. Esempio semplificato:

\begin{quote}
    TB $\leq$ 1.2? \\
    $\quad \downarrow$ sì \\
    $\quad$ controlla Sgpt \\
    $\quad \quad \downarrow$ Sgpt $\leq$ 35 \\
    $\quad \quad \downarrow$ sì \\
    $\quad \quad$ predici classe 1 \\
    $\quad \quad \downarrow$ no \\
    $\quad \quad$ predici classe 0 \\
    $\quad \downarrow$ no \\
    $\quad$ controlla DB \\
    $\quad \quad \downarrow$ DB $\leq$ 0.5 \\
    $\quad \quad \downarrow$ sì \\
    $\quad \quad$ predici classe 0 \\
    $\quad \quad \downarrow$ no \\
    $\quad \quad$ predici classe 0
\end{quote}

Ogni nodo interno dell'albero contiene una condizione. Ogni foglia contiene una predizione.

\section{Come decide un albero decisionale?}
Un albero decisionale cerca di dividere i dati in gruppi sempre più ``puri''. Un gruppo è puro quando contiene quasi solo esempi della stessa classe.

Per misurare l'impurità, \texttt{DecisionTreeClassifier} usa di default l'\textit{indice di Gini}. La formula è:

\begin{equation}
    \text{Gini} = 1 - \sum_{i=1}^{n} p_i^2
\end{equation}

dove $p_i$ è la proporzione della classe $i$ nel nodo.

\begin{example}
Supponiamo un nodo con:
\begin{align*}
    50 \text{ pazienti classe } 0 \\
    50 \text{ pazienti classe } 1
\end{align*}

Allora:
\begin{align*}
    p_0 &= 0.5 \\
    p_1 &= 0.5 \\
    \text{Gini} &= 1 - (0.5^2 + 0.5^2) \\
               &= 1 - (0.25 + 0.25) \\
               &= 0.5
\end{align*}

Il nodo è molto impuro.
\end{example}

\begin{example}
Supponiamo invece:
\begin{align*}
    100 \text{ pazienti classe } 0 \\
    0 \text{ pazienti classe } 1
\end{align*}

Allora:
\begin{align*}
    p_0 &= 1 \\
    p_1 &= 0 \\
    \text{Gini} &= 1 - (1^2 + 0^2) \\
               &= 0
\end{align*}

Il nodo è perfettamente puro.
\end{example}

L'albero cerca gli \textit{split} che riducono il più possibile l'impurità.

\section{Valutazione del primo albero}
Dopo l'addestramento, il modello viene valutato sul test set:

\begin{lstlisting}[language=Python]
from sklearn.metrics import confusion_matrix

y_pred = model.predict(X_test)
cm = confusion_matrix(y_test, y_pred)

print("Confusion matrix:\n", cm)
print("\nAccuracy:", cm.diagonal().sum() / cm.sum())
\end{lstlisting}

La matrice di confusione ottenuta è:
\begin{align*}
    \begin{bmatrix}
        96 & 29 \\
        28 & 22
    \end{bmatrix}
\end{align*}

Questa matrice va letta così:

\begin{table}[H]
\centering
\caption{Matrice di confusione del primo albero decisionale}
\label{tab:confusion_matrix_1}
\begin{tabular}{lcc}
\toprule
 & \textbf{Predetto 0} & \textbf{Predetto 1} \\
\midrule
\textbf{Vero 0} & 96 & 29 \\
\textbf{Vero 1} & 28 & 22 \\
\bottomrule
\end{tabular}
\end{table}

Dato che:
\begin{align*}
    0 &= \text{Liver Disease} \\
    1 &= \text{No Liver Disease}
\end{align*}

abbiamo:
\begin{itemize}
    \item 96 pazienti con \texttt{Liver Disease} classificati correttamente
    \item 29 pazienti con \texttt{Liver Disease} classificati erroneamente come \texttt{No Liver Disease}
    \item 28 pazienti con \texttt{No Liver Disease} classificati erroneamente come \texttt{Liver Disease}
    \item 22 pazienti con \texttt{No Liver Disease} classificati correttamente
\end{itemize}

L'\textit{accuracy} è:
\begin{align*}
    \text{accuracy} &\approx 0.6743 \text{ (circa 67.4\%)}
\end{align*}

Il risultato non è ottimo. Il modello riconosce abbastanza la classe maggioritaria, ma fatica con la classe minoritaria.

\section{Accuracy, precision, recall e F1-score}
\subsection{Accuracy}
L'\textit{accuracy} misura quante predizioni sono corrette sul totale:
\begin{equation}
    \text{accuracy} = \frac{\text{predizioni corrette}}{\text{predizioni totali}}
\end{equation}

Nel primo modello:
\begin{align*}
    \text{accuracy} &= \frac{96 + 22}{175} \\
                    &= \frac{118}{175} \\
                    &\approx 0.674
\end{align*}

\subsection{Precision}
La \textit{precision} risponde alla domanda:
\begin{quote}
    tra tutti quelli che il modello ha predetto come classe positiva, quanti erano davvero positivi?
\end{quote}

Per la classe 1:
\begin{equation}
    \text{precision}_1 = \frac{\text{veri 1 predetti 1}}{\text{tutti i predetti 1}}
\end{equation}

Nel primo modello:
\begin{align*}
    \text{precision}_1 &= \frac{22}{29 + 22} \\
                       &= \frac{22}{51} \\
                       &\approx 0.43
\end{align*}

\subsection{Recall}
La \textit{recall} risponde alla domanda:
\begin{quote}
    tra tutti quelli che erano davvero classe positiva, quanti ne ha trovati il modello?
\end{quote}

Per la classe 1:
\begin{equation}
    \text{recall}_1 = \frac{\text{veri 1 predetti 1}}{\text{tutti i veri 1}}
\end{equation}

Nel primo modello:
\begin{align*}
    \text{recall}_1 &= \frac{22}{28 + 22} \\
                    &= \frac{22}{50} \\
                    &= 0.44
\end{align*}

\subsection{F1-score}
L'\textit{F1-score} è una media armonica tra precision e recall:
\begin{equation}
    F1 = 2 \cdot \frac{\text{precision} \cdot \text{recall}}{\text{precision} + \text{recall}}
\end{equation}

È utile quando vogliamo un compromesso tra precision e recall.

\section{Problema del primo albero: overfitting}
Un albero decisionale senza vincoli può crescere troppo. Se cresce troppo, può imparare quasi a memoria il training set. Questo fenomeno si chiama \textit{overfitting}.

Un modello in \textit{overfitting}:
\begin{itemize}
    \item va molto bene sul training set
    \item va male o non migliora sul test set
\end{itemize}

Nel notebook viene visualizzato l'albero con:

\begin{lstlisting}[language=Python]
from sklearn.tree import plot_tree
import matplotlib.pyplot as plt

fig, ax = plt.subplots(figsize=(150, 100))
plot_tree(model, filled=True, ax=ax)
plt.plot()
\end{lstlisting}

La dimensione enorme della figura suggerisce che l'albero completo sia troppo grande e difficile da interpretare.

\section{Miglioramento: limitare la profondità dell'albero}
Per controllare l'\textit{overfitting} si prova un albero con profondità massima:

\begin{lstlisting}[language=Python]
prunedModel = DecisionTreeClassifier(
    random_state=42,
    max_depth=5,
    min_samples_split=2
)
\end{lstlisting}

Il parametro:
\begin{lstlisting}[language=Python]
max_depth=5
\end{lstlisting}

significa:
\begin{quote}
    l'albero può avere al massimo profondità 5
\end{quote}

Questo rende il modello più semplice. Tuttavia il risultato ottenuto è:
\begin{align*}
    \text{Confusion matrix:} &\begin{bmatrix} 104 & 21 \\ 41 & 9 \end{bmatrix} \\
    \text{Accuracy:} &\approx 0.6457
\end{align*}

Quindi l'\textit{accuracy} scende a circa 64.6\%. La classe 1 viene riconosciuta molto male:
\begin{align*}
    \text{solo 9 veri positivi su 50}
\end{align*}

Questo modello è troppo sbilanciato verso la classe 0.

\section{Pruning tramite \texttt{ccp\_alpha}}
Il notebook usa poi il parametro di complessità \texttt{ccp\_alpha}:

\begin{lstlisting}[language=Python]
path = model.cost_complexity_pruning_path(X_train, y_train)
ccp_alphas = path.ccp_alphas
print(ccp_alphas)
\end{lstlisting}

Il \textit{pruning} serve a potare l'albero, cioè a eliminare rami poco utili. L'idea è:
\begin{quote}
    un albero più grande può essere più preciso sul training set, ma può generalizzare peggio.
\end{quote}

Il parametro \texttt{ccp\_alpha} controlla quanto penalizziamo la complessità dell'albero. La funzione obiettivo diventa, in modo semplificato:
\begin{equation}
    \text{errore dell'albero} + \alpha \cdot \text{complessità dell'albero}
\end{equation}

Se $\alpha$ è piccolo, l'albero può rimanere grande. Se $\alpha$ è grande, l'albero viene potato di più.

\section{Ricerca degli iperparametri}
Nel notebook viene fatta una ricerca manuale su più valori di:
\begin{itemize}
    \item \texttt{max\_depth}
    \item \texttt{ccp\_alpha}
    \item \texttt{class\_weight}
\end{itemize}

Per esempio:

\begin{lstlisting}[language=Python]
for depth in range(4, 8):
    for complexity in ccp_alphas:
        clf = DecisionTreeClassifier(
            max_depth=depth,
            ccp_alpha=complexity,
            class_weight="balanced",
            random_state=42
        )
\end{lstlisting}

Questa parte è una forma di \textit{grid search} manuale. Si provano molte configurazioni e si sceglie quella che ottiene il risultato migliore sul test set.

Il primo risultato migliore trovato è:
\begin{align*}
    \text{Miglior max\_depth:} &= 6 \\
    \text{Miglior ccp\_alpha:} &= 0.0 \\
    \text{Train accuracy:} &= 0.7745 \\
    \text{Test accuracy:} &= 0.6914
\end{align*}

La matrice di confusione è:
\begin{align*}
    \begin{bmatrix}
        79 & 46 \\
        8 & 42
    \end{bmatrix}
\end{align*}

Rispetto al primo albero:
\begin{quote}
    classe 1 riconosciuta molto meglio
\end{quote}

Infatti:
\begin{align*}
    \text{veri classe 1} &= 50 \\
    \text{classe 1 correttamente predetti} &= 42
\end{align*}

La \textit{recall} della classe 1 diventa:
\begin{align*}
    \text{recall}_1 &= \frac{42}{50} = 0.84
\end{align*}

Quindi il modello trova l'84\% dei casi della classe 1. Però aumenta il numero di falsi positivi:
\begin{align*}
    46 \text{ pazienti classe 0 vengono predetti come classe 1}
\end{align*}

Questo è il classico compromesso:
\begin{quote}
    aumento la recall della classe minoritaria, ma rischio più errori sulla classe maggioritaria.
\end{quote}

\section{Balanced accuracy}
Nel notebook viene calcolata anche la \textit{balanced accuracy}:

\begin{lstlisting}[language=Python]
from sklearn.metrics import balanced_accuracy_score

print("Balanced accuracy:", balanced_accuracy_score(y_test, y_pred))
\end{lstlisting}

La \textit{balanced accuracy} è la media delle recall delle classi. Per due classi:
\begin{equation}
    \text{balanced accuracy} = \frac{\text{recall classe 0} + \text{recall classe 1}}{2}
\end{equation}

Con la matrice:
\begin{align*}
    \begin{bmatrix}
        79 & 46 \\
        8 & 42
    \end{bmatrix}
\end{align*}

abbiamo:
\begin{align*}
    \text{recall classe 0} &= \frac{79}{79 + 46} = \frac{79}{125} = 0.632 \\
    \text{recall classe 1} &= \frac{42}{8 + 42} = \frac{42}{50} = 0.84
\end{align*}

Quindi:
\begin{align*}
    \text{balanced accuracy} &= \frac{0.632 + 0.84}{2} \\
                              &= 0.736
\end{align*}

Il notebook infatti ottiene:
\begin{quote}
    Balanced accuracy: 0.736
\end{quote}

Questa metrica è più adatta dell'\textit{accuracy} semplice quando le classi sono sbilanciate.

\section{Uso dei pesi di classe}
Il progetto prova poi a dare più importanza alla classe minoritaria usando \texttt{class\_weight}. Nel notebook viene definita una funzione:

\begin{lstlisting}[language=Python]
def train_best_weighted_tree(...):
\end{lstlisting}

La funzione prova diversi pesi:
\begin{lstlisting}[language=Python]
weights_to_try = [
    {0: 1, 1: 1},
    {0: 1, 1: 1.5},
    {0: 1, 1: 2},
    {0: 1, 1: 2.5},
    {0: 1, 1: 3},
    {0: 1, 1: 4},
    {0: 1, 1: 4.5},
    {0: 1, 1: 5},
    {0: 1, 1: 7},
    {0: 1, 1: 9},
    {0: 1, 1: 10}
]
\end{lstlisting}

Un peso:
\begin{lstlisting}[language=Python]
{0: 1, 1: 2.5}
\end{lstlisting}

significa:
\begin{quote}
    gli errori sulla classe 1 pesano 2.5 volte di più degli errori sulla classe 0
\end{quote}

Questo spinge l'albero a prestare più attenzione alla classe 1.

\section{Miglior albero pesato}
La migliore configurazione trovata è:
\begin{align*}
    \text{max\_depth:} &= 6 \\
    \text{ccp\_alpha:} &= 0.0 \\
    \text{class\_weight:} &= \{0: 1, 1: 2.5\} \\
    \text{train\_accuracy:} &= 0.7745 \\
    \text{test\_accuracy:} &= 0.6914 \\
    \text{test\_balanced\_accuracy:} &= 0.736 \\
    \text{precision\_class\_1:} &= 0.4773 \\
    \text{recall\_class\_1:} &= 0.84 \\
    \text{f1\_class\_1:} &= 0.6087 \\
    \text{macro\_f1:} &= 0.6770
\end{align*}

Matrice di confusione:
\begin{align*}
    \begin{bmatrix}
        79 & 46 \\
        8 & 42
    \end{bmatrix}
\end{align*}

Il risultato è lo stesso già visto con \texttt{class\_weight="balanced"}. Questo modello è migliore del primo albero semplice se ci interessa riconoscere bene la classe 1.

Però bisogna stare attenti: la classe 1 è \texttt{No Liver Disease}. Se l'obiettivo medico fosse riconoscere soprattutto i malati, probabilmente bisognerebbe invertire la codifica o scegliere come positiva la classe \texttt{Liver Disease}.

\section{Ensemble di tre alberi con majority voting}
Il notebook prova poi una tecnica di \textit{ensemble}. Un \textit{ensemble} combina più modelli invece di usarne uno solo. Qui vengono addestrati tre alberi decisionali e le loro predizioni vengono combinate tramite voto di maggioranza.

La funzione principale è:
\begin{lstlisting}[language=Python]
def majority_vote(predictions):
\end{lstlisting}

L'idea è semplice. Supponiamo che tre modelli predicano:
\begin{align*}
    \text{modello 1} &\rightarrow 1 \\
    \text{modello 2} &\rightarrow 0 \\
    \text{modello 3} &\rightarrow 1
\end{align*}

Due modelli su tre hanno votato \texttt{1}, quindi la predizione finale è:
\begin{align*}
    1
\end{align*}

Se invece:
\begin{align*}
    \text{modello 1} &\rightarrow 0 \\
    \text{modello 2} &\rightarrow 0 \\
    \text{modello 3} &\rightarrow 1
\end{align*}

la predizione finale è:
\begin{align*}
    0
\end{align*}

Nel codice:
\begin{lstlisting}[language=Python]
if votes_for_1 >= 2:
    final_predictions.append(1)
else:
    final_predictions.append(0)
\end{lstlisting}

\section{Addestramento dei tre alberi}
La funzione:
\begin{lstlisting}[language=Python]
def train_three_trees(...)
\end{lstlisting}

usa:
\begin{lstlisting}[language=Python]
from sklearn.model_selection import StratifiedKFold

StratifiedKFold(n_splits=3, shuffle=True, random_state=random_state)
\end{lstlisting}

Quindi il training set viene diviso in 3 parti stratificate. Ogni albero viene addestrato su una parte diversa dei dati. Questo crea tre modelli leggermente diversi.

L'idea generale è:
\begin{quote}
    un singolo albero può essere instabile; più alberi possono compensarsi tra loro.
\end{quote}

Questo principio è alla base anche delle \textit{Random Forest}.

\section{Risultato dell'ensemble}
La migliore configurazione trovata senza \textit{undersampling} è:
\begin{align*}
    \text{max\_depth:} &= 5 \\
    \text{ccp\_alpha:} &= 0.0 \\
    \text{class\_weight:} &= \{0: 1, 1: 5\} \\
    \text{test\_accuracy:} &\approx 0.6057 \\
    \text{test\_balanced\_accuracy:} &\approx 0.706 \\
    \text{precision\_class\_1:} &\approx 0.4159 \\
    \text{recall\_class\_1:} &= 0.94 \\
    \text{f1\_class\_1:} &\approx 0.5767
\end{align*}

Matrice di confusione:
\begin{align*}
    \begin{bmatrix}
        59 & 66 \\
        3 & 47
    \end{bmatrix}
\end{align*}

Interpretazione:
\begin{align*}
    47 \text{ elementi della classe 1 su 50 vengono riconosciuti correttamente}
\end{align*}

La \textit{recall} della classe 1 è altissima:
\begin{align*}
    \text{recall}_1 &= \frac{47}{50} = 0.94
\end{align*}

Però il modello predice troppi esempi come classe 1. Infatti:
\begin{align*}
    66 \text{ esempi della classe 0 vengono classificati erroneamente come classe 1}
\end{align*}

Quindi l'\textit{ensemble} aumenta molto la sensibilità verso la classe 1, ma peggiora la \textit{precision}.

\section{Undersampling}
Il notebook prova poi una tecnica diversa per gestire lo sbilanciamento: l'\textit{undersampling}. L'\textit{undersampling} consiste nel ridurre il numero di esempi della classe maggioritaria, in modo da avere un training set bilanciato.

La funzione è:
\begin{lstlisting}[language=Python]
def train_test_split_with_undersampling(...)
\end{lstlisting}

La cosa corretta fatta nel notebook è questa:
\begin{quote}
    prima si divide in train/test, poi si applica \textit{undersampling} solo sul training set.
\end{quote}

Questo è molto importante. Non bisogna fare \textit{undersampling} prima dello split, perché si rischia di alterare artificialmente anche il test set. Il test set deve rimanere realistico.

\section{Come funziona l'undersampling nel codice}
Il codice fa:

\begin{lstlisting}[language=Python]
class_counts = train_df[target_col].value_counts()
min_class_count = class_counts.min()
\end{lstlisting}

Trova la classe meno numerosa. Poi, per ogni classe, prende lo stesso numero di esempi:

\begin{lstlisting}[language=Python]
class_df_under = class_df.sample(
    n=min_class_count,
    random_state=random_state
)
\end{lstlisting}

Alla fine il training set bilanciato contiene:
\begin{align*}
    0 &\rightarrow 120 \\
    1 &\rightarrow 120
\end{align*}

Quindi il modello viene addestrato su:
\begin{align*}
    120 \text{ esempi classe 0} \\
    120 \text{ esempi classe 1}
\end{align*}

\section{Ensemble dopo undersampling}
Dopo l'\textit{undersampling} si riesegue la ricerca dell'\textit{ensemble}. Il risultato migliore è:
\begin{align*}
    \text{max\_depth:} &= 5 \\
    \text{ccp\_alpha:} &= 0.0 \\
    \text{class\_weight:} &= \{0: 1, 1: 1.5\} \\
    \text{test\_accuracy:} &\approx 0.6914 \\
    \text{test\_balanced\_accuracy:} &\approx 0.7756 \\
    \text{precision\_class\_1:} &\approx 0.4639 \\
    \text{recall\_class\_1:} &\approx 0.9574 \\
    \text{f1\_class\_1:} &\approx 0.625 \\
    \text{macro\_f1:} &\approx 0.6814 \\
    \text{weighted\_f1:} &\approx 0.7076
\end{align*}

Matrice di confusione:
\begin{align*}
    \begin{bmatrix}
        76 & 52 \\
        2 & 45
    \end{bmatrix}
\end{align*}

Questo è uno dei risultati più interessanti del notebook. La classe 1 viene riconosciuta quasi sempre:
\begin{align*}
    \text{recall}_1 &\approx \frac{45}{47} = 0.957
\end{align*}

Quindi la \textit{recall} della classe 1 è circa 95.7\%. La \textit{balanced accuracy} è:
\begin{align*}
    0.7756
\end{align*}

che è migliore rispetto ai modelli precedenti. Anche qui, però, la \textit{precision} della classe 1 resta bassa:
\begin{align*}
    0.4639
\end{align*}

cioè molti pazienti vengono predetti come classe 1 anche quando non lo sono.

\section{Random Forest}
Il notebook prova infine una \textit{Random Forest}. Una \textit{Random Forest} è un insieme di molti alberi decisionali. Mentre prima si usavano tre alberi, qui si usano molti più alberi, per esempio:

\begin{lstlisting}[language=Python]
n_estimators_to_try = [100, 200, 300, 500]
\end{lstlisting}

Ogni albero viene addestrato con una certa casualità. La previsione finale è ottenuta combinando le predizioni di tutti gli alberi. In classificazione, di solito, si usa il voto di maggioranza.

\section{Perché Random Forest può funzionare meglio di un singolo albero?}
Un singolo albero decisionale è molto sensibile ai dati. Piccoli cambiamenti nel training set possono produrre un albero molto diverso. La \textit{Random Forest} riduce questo problema perché combina molti alberi.

L'idea è:
\begin{quote}
    un albero singolo può sbagliare in modo instabile; molti alberi insieme producono una previsione più robusta.
\end{quote}

In generale:
\begin{align*}
    \text{Decision Tree} &= \text{modello semplice ma instabile} \\
    \text{Random Forest} &= \text{insieme di alberi, più stabile}
\end{align*}

\section{Ricerca degli iperparametri della Random Forest}
Nel notebook viene definita la funzione:
\begin{lstlisting}[language=Python]
def search_best_random_forest(...)
\end{lstlisting}

La funzione prova combinazioni di:
\begin{itemize}
    \item \texttt{n\_estimators}
    \item \texttt{max\_depth}
    \item \texttt{class\_weight}
\end{itemize}

I valori provati sono:
\begin{lstlisting}[language=Python]
n_estimators_to_try = [100, 200, 300, 500]
depths_to_try = [4, 5, 6, 7, 8, None]
weights_to_try = [
    {0: 1, 1: 1},
    {0: 1, 1: 1.5},
    {0: 1, 1: 2},
    {0: 1, 1: 3},
    {0: 1, 1: 5},
    "balanced"
]
\end{lstlisting}

La metrica usata per scegliere il modello migliore è:
\begin{lstlisting}[language=Python]
scoring_metric="macro_f1"
\end{lstlisting}

Il \textit{macro-F1} è la media degli F1-score delle classi, senza pesare di più la classe più frequente. È utile quando vogliamo valutare il modello in modo più equo su entrambe le classi.

\section{Risultato della Random Forest}
La migliore configurazione trovata è:
\begin{align*}
    \text{n\_estimators:} &= 200 \\
    \text{max\_depth:} &= \text{None} \\
    \text{class\_weight:} &= \{0: 1, 1: 1\} \\
    \text{train\_accuracy:} &= 1.0 \\
    \text{test\_accuracy:} &\approx 0.6857 \\
    \text{test\_balanced\_accuracy:} &\approx 0.7178 \\
    \text{precision\_class\_1:} &\approx 0.4512 \\
    \text{recall\_class\_1:} &\approx 0.7872 \\
    \text{f1\_class\_1:} &\approx 0.5736 \\
    \text{macro\_f1:} &\approx 0.6624
\end{align*}

Matrice di confusione:
\begin{align*}
    \begin{bmatrix}
        83 & 45 \\
        10 & 37
    \end{bmatrix}
\end{align*}

Interpretazione:
\begin{align*}
    83 \text{ elementi classe 0 corretti} \\
    45 \text{ elementi classe 0 sbagliati} \\
    10 \text{ elementi classe 1 sbagliati} \\
    37 \text{ elementi classe 1 corretti}
\end{align*}

La \textit{Random Forest} ottiene una \textit{accuracy} simile agli altri modelli, ma non supera l'\textit{ensemble} con \textit{undersampling} sulla \textit{balanced accuracy}. Inoltre:
\begin{align*}
    \text{train\_accuracy} &= 1.0
\end{align*}

Questo indica possibile \textit{overfitting}. Il modello impara perfettamente il training set, ma sul test set arriva solo a circa 68.6\%.

\section{Confronto tra i modelli principali}
\begin{table}[H]
\centering
\caption{Confronto tra i modelli principali}
\label{tab:model_comparison}
\begin{tabular}{>{\raggedright\arraybackslash}p{4cm}ccccc>{\raggedright\arraybackslash}p{4cm}}
\toprule
\textbf{Modello} & \textbf{Accuracy test} & \textbf{Balanced accuracy} & \textbf{Recall classe 1} & \textbf{F1 classe 1} & \textbf{Nota} \\
\midrule
Albero base & 0.674 & 0.604 & 0.44 & 0.44 & modello semplice, classe 1 debole \\
Albero max\_depth=5 & 0.646 & 0.486 & 0.18 & 0.23 & troppo orientato verso classe 0 \\
Albero pesato & 0.691 & 0.736 & 0.84 & 0.609 & buon compromesso \\
Ensemble 3 alberi & 0.606 & 0.706 & 0.94 & 0.577 & recall altissima, precision bassa \\
Ensemble + undersampling & 0.691 & 0.776 & 0.957 & 0.625 & miglior risultato sulla classe 1 \\
Random Forest & 0.686 & 0.718 & 0.787 & 0.574 & stabile ma possibile overfitting \\
\bottomrule
\end{tabular}
\end{table}

Il modello più interessante nel notebook è:
\begin{quote}
    \textbf{Ensemble di tre alberi + undersampling}
\end{quote}

perché ottiene:
\begin{align*}
    \text{balanced accuracy} &\approx 0.7756 \\
    \text{recall classe 1} &\approx 0.9574 \\
    \text{f1 classe 1} &\approx 0.625
\end{align*}

È il modello che riconosce meglio la classe 1.

\section{Feature importance con Mutual Information}
Alla fine il notebook calcola l'importanza delle variabili usando la \textit{mutual information}:

\begin{lstlisting}[language=Python]
from sklearn.feature_selection import mutual_info_classif
\end{lstlisting}

La \textit{mutual information} misura quanta informazione una variabile fornisce sul target. In modo intuitivo:
\begin{quote}
    se conoscere una variabile aiuta molto a prevedere il target, allora quella variabile ha \textit{mutual information} alta.
\end{quote}

Il risultato ottenuto è:
\begin{align*}
    \text{Sgpt}        &\rightarrow 0.072161 \\
    \text{Sgot}        &\rightarrow 0.068853 \\
    \text{Age}         &\rightarrow 0.066943 \\
    \text{TB}          &\rightarrow 0.053878 \\
    \text{DB}          &\rightarrow 0.048560 \\
    \text{Alkphos}     &\rightarrow 0.028152 \\
    \text{ALB}         &\rightarrow 0.025456 \\
    \text{A/G Ratio}   &\rightarrow 0.018387 \\
    \text{Gender}      &\rightarrow 0.000000 \\
    \text{TP}          &\rightarrow 0.000000
\end{align*}

Quindi le variabili più informative risultano:
\begin{itemize}
    \item \texttt{Sgpt}
    \item \texttt{Sgot}
    \item \texttt{Age}
    \item \texttt{TB}
    \item \texttt{DB}
\end{itemize}

Questa cosa ha anche senso dal punto di vista clinico, perché SGPT, SGOT e bilirubina sono misure legate al funzionamento del fegato.

\section{Attenzione a un piccolo errore nel codice della mutual information}
Nel notebook c'è questa parte:
\begin{lstlisting}[language=Python]
categoriche = ["gender"]
\end{lstlisting}

Però la colonna si chiama:
\begin{lstlisting}[language=Python]
"Gender"
\end{lstlisting}

con la G maiuscola. Quindi sarebbe meglio scrivere:
\begin{lstlisting}[language=Python]
categoriche = ["Gender"]
\end{lstlisting}

Versione corretta:
\begin{lstlisting}[language=Python]
categoriche = ["Gender"]

discrete_mask = []

for col in X.columns:
    if col in categoriche:
        discrete_mask.append(True)
    else:
        discrete_mask.append(False)
\end{lstlisting}

Dato che \texttt{Gender} è binaria, l'errore non distrugge completamente l'analisi, ma è comunque meglio correggerlo.

\section{Sporcamento del dataset}
Il notebook introduce anche una funzione per ``sporcare'' variabili categoriche:
\begin{lstlisting}[language=Python]
def corrupt_rows_cat(df, col, frac=0.1, random_state=None):
\end{lstlisting}

L'idea è simulare rumore nei dati. Per esempio, se \texttt{frac=0.1}, si modifica casualmente il 10\% delle righe di una certa colonna. Nel caso di una variabile categorica come \texttt{Gender}, il codice può cambiare alcuni valori:
\begin{align*}
    0 &\rightarrow 1 \\
    1 &\rightarrow 0
\end{align*}

Questo serve per studiare quanto il modello è robusto quando i dati contengono errori.

Esempio pratico:
\begin{lstlisting}[language=Python]
df_sporco = corrupt_rows_cat(
    df=df_encoded,
    col="Gender",
    frac=0.1,
    random_state=42
)
\end{lstlisting}

Questo crea una copia del dataset in cui il 10\% dei valori della colonna \texttt{Gender} viene modificato.

\section{Attenzione a un errore nella funzione \texttt{corrupt\_rows\_cat}}
Nel codice c'è questa riga:
\begin{lstlisting}[language=Python]
if random_state != None: random_state = None
\end{lstlisting}

Questa riga annulla il \texttt{random\_state}. Probabilmente l'intenzione era rendere la funzione riproducibile, ma così facendo si ottiene l'effetto contrario.

Una versione più corretta sarebbe:
\begin{lstlisting}[language=Python]
import numpy as np

def corrupt_rows_cat(df, col, frac=0.1, random_state=None):
    corrupt_df = df.copy()
    rng = np.random.default_rng(random_state)
    unique_values = corrupt_df[col].unique()

    if len(unique_values) <= 1:
        print(f"Attenzione: la colonna '{col}' ha un solo valore unico.")
        return corrupt_df

    n_rows_to_corrupt = int(frac * len(corrupt_df))
    indices_to_corrupt = rng.choice(
        corrupt_df.index,
        size=n_rows_to_corrupt,
        replace=False
    )

    for idx in indices_to_corrupt:
        current_value = corrupt_df.at[idx, col]
        pool_of_new_value = unique_values[unique_values != current_value]
        corrupt_df.at[idx, col] = rng.choice(pool_of_new_value)

    return corrupt_df
\end{lstlisting}

Così, usando lo stesso \texttt{random\_state}, si ottiene sempre lo stesso dataset corrotto.

\section{Osservazione importante: quale classe stiamo ottimizzando?}
Questo è probabilmente il punto più importante da spiegare in una relazione.

Nel notebook:
\begin{align*}
    0 &= \text{Liver Disease} \\
    1 &= \text{No Liver Disease}
\end{align*}

Però molte funzioni ottimizzano:
\begin{lstlisting}[language=Python]
scoring_metric="f1_class_1"
\end{lstlisting}

Quindi stanno ottimizzando la classe:
\begin{quote}
    \texttt{No Liver Disease}
\end{quote}

Questo non è sbagliato in assoluto, ma va dichiarato. In un contesto medico, spesso ci interessa soprattutto non perdere i malati. In quel caso potrebbe essere più naturale codificare:
\begin{align*}
    \text{Liver Disease}     &\rightarrow 1 \\
    \text{No Liver Disease}  &\rightarrow 0
\end{align*}

Così la classe positiva sarebbe la malattia. Una versione più chiara sarebbe:
\begin{lstlisting}[language=Python]
df_encoded["Selector"] = df["Selector"].map({
    "No Liver Disease": 0,
    "Liver Disease": 1
})
\end{lstlisting}

In questo modo:
\begin{align*}
    1 &= \text{paziente malato} \\
    0 &= \text{paziente non malato}
\end{align*}

E allora metriche come \texttt{recall\_class\_1} avrebbero un significato clinico più naturale:
\begin{quote}
    quanti pazienti malati riesco a individuare?
\end{quote}

\section{Nota su \texttt{classification\_report}}
In alcune celle del notebook compare:
\begin{lstlisting}[language=Python]
print(classification_report(y_pred, y_test))
\end{lstlisting}

L'ordine corretto normalmente è:
\begin{lstlisting}[language=Python]
print(classification_report(y_test, y_pred))
\end{lstlisting}

La funzione si usa così:
\begin{lstlisting}[language=Python]
classification_report(y_true, y_pred)
\end{lstlisting}

dove:
\begin{align*}
    y\_true &= \text{valori reali} \\
    y\_pred &= \text{predizioni del modello}
\end{align*}

Quindi la forma corretta è:
\begin{lstlisting}[language=Python]
print(classification_report(y_test, y_pred))
\end{lstlisting}

Nel notebook in alcune celle l'ordine viene poi usato correttamente.

\section{Struttura generale del progetto}
Possiamo riassumere il progetto come una \textit{pipeline} di \textit{machine learning}:
\begin{enumerate}
    \item Importazione librerie
    \item Download e lettura dataset
    \item Esplorazione iniziale dei dati
    \item Analisi della distribuzione del target
    \item Codifica delle variabili categoriche
    \item Train/test split
    \item Addestramento di un Decision Tree base
    \item Valutazione con confusion matrix e metriche
    \item Miglioramento tramite max\_depth e pruning
    \item Ricerca di iperparametri
    \item Gestione dello sbilanciamento con class\_weight
    \item Ensemble di tre alberi
    \item Undersampling del training set
    \item Random Forest
    \item Analisi dell'importanza delle variabili
    \item Simulazione di rumore nei dati
\end{enumerate}

\section{Conclusione}
Il progetto mostra un percorso abbastanza completo di classificazione binaria su dati clinici.

Il punto di partenza è un semplice albero decisionale, che ottiene circa il 67\% di \textit{accuracy}, ma fatica a riconoscere bene la classe minoritaria.

Successivamente vengono introdotte tecniche più avanzate:
\begin{itemize}
    \item pruning
    \item class\_weight
    \item balanced accuracy
    \item ricerca degli iperparametri
    \item ensemble
    \item undersampling
    \item random forest
    \item mutual information
\end{itemize}

Il risultato migliore, considerando la \textit{balanced accuracy} e la capacità di riconoscere la classe 1, viene ottenuto con:
\begin{quote}
    \textbf{ensemble di tre alberi + undersampling}
\end{quote}

che raggiunge:
\begin{align*}
    \text{accuracy} &\approx 0.691 \\
    \text{balanced accuracy} &\approx 0.776 \\
    \text{recall classe 1} &\approx 0.957 \\
    \text{f1 classe 1} &\approx 0.625
\end{align*}

Tuttavia bisogna interpretare bene la codifica delle classi:
\begin{align*}
    \text{classe 0} &= \text{Liver Disease} \\
    \text{classe 1} &= \text{No Liver Disease}
\end{align*}

Quindi il progetto, così com'è scritto, ottimizza spesso il riconoscimento della classe \texttt{No Liver Disease}. Per renderlo più coerente con un problema medico, sarebbe consigliabile ricodificare il target in modo che:
\begin{align*}
    1 &= \text{Liver Disease} \\
    0 &= \text{No Liver Disease}
\end{align*}

e poi ottimizzare la \textit{recall} o l'\textit{F1} della classe malata.

In sintesi, il progetto è buono perché non si limita ad addestrare un singolo modello, ma esplora diversi problemi reali del \textit{machine learning}:
\begin{itemize}
    \item dataset sbilanciato
    \item overfitting
    \item scelta delle metriche
    \item interpretabilità
    \item ricerca degli iperparametri
    \item robustezza ai dati rumorosi
\end{itemize}

La parte più importante da valorizzare nella relazione è che \textbf{l'\textit{accuracy} da sola non basta}: in un problema clinico bisogna guardare attentamente \textit{precision}, \textit{recall}, \textit{F1-score}, \textit{balanced accuracy} e matrice di confusione.

\end{document}
```